\chapter{Acne}

\section*{Definition}
Acne vulgaris is a common inflammatory disorder of the pilosebaceous units. It causes comedones (blackheads and whiteheads), papules, pustules, and, in more severe disease, nodules. It most often affects the face, chest, and back and may lead to scarring or persistent color changes.

\section*{When to See a Doctor}

\begin{itemize}
 \item Acne that is severe, painful, nodular, scarring, widespread, or causing significant distress
 \item Acne that does not improve after an adequate trial of appropriate treatment, or that repeatedly returns
 \item Sudden or unusually severe acne in an adult, or acne with irregular periods, excess facial or body hair, hair thinning, or other signs of androgen excess
 \item Acne associated with fever, joint pain, or rapidly worsening painful lesions
\end{itemize}

\section*{How It Develops}
Acne develops when excess sebum, abnormal shedding of cells lining the follicle, microbial changes within the follicle, and inflammation combine to block pores and produce lesions. Androgens increase sebaceous-gland activity, particularly during puberty. Adult acne can occur with hormonal changes, medicines, or no identifiable hormonal disorder; perimenopause and menopause may alter the balance between androgen and estrogen effects but do not have one single mechanism.

\section*{Causes}
\begin{itemize}
 \item Hormonal changes during puberty, the menstrual cycle, pregnancy, perimenopause, or menopause; PCOS may contribute in some people
 \item Increased androgen activity or an endocrine disorder causing androgen excess
 \item Genetic predisposition and increased sebum production
 \item Oily or comedogenic cosmetics, hair products, or skin-care products
 \item Medicines such as systemic corticosteroids, lithium, androgens, and some progestin-containing contraceptives; contraceptive effects vary by formulation
 \item Mechanical occlusion or friction from helmets, headbands, masks, tight clothing, or equipment
 \item Stress may worsen existing acne in some people but is not the sole cause
 \item A high-glycemic diet or frequent intake of certain dairy products may worsen acne in some people, but dietary effects vary and evidence is not definitive
 \item Harsh scrubbing, picking, or over-washing can irritate skin and worsen inflammation but does not cause acne by itself
\end{itemize}

\section*{Related Symptoms}
\begin{itemize}
 \item Oily skin, blackheads, whiteheads, papules, pustules, nodules, or tenderness
 \item Post-inflammatory dark or red marks and acne scars
 \item Hirsutism, scalp hair thinning, or irregular periods, which may suggest androgen excess or PCOS rather than being symptoms of acne itself
 \item Embarrassment, low mood, anxiety, or reduced quality of life
\end{itemize}


\section*{Medical Evaluation}
\begin{itemize}
 \item Acne is usually diagnosed clinically by examining the distribution and type of lesions, severity, scarring, skin type, medicines, and possible acneiform disorders.
 \item Hormone testing is not routine. Testosterone, DHEAS, or other endocrine tests may be considered when acne is sudden, severe, persistent in adulthood, or accompanied by irregular periods, hirsutism, virilization, or other signs of androgen excess.
 \item A bacterial swab or culture is not routinely useful for ordinary acne; testing may be considered when an alternative diagnosis, unusual folliculitis, or secondary infection is suspected.
 \item Referral to a dermatologist is appropriate for severe, scarring, psychologically distressing, atypical, or treatment-resistant acne and when isotretinoin or procedural scar treatment is being considered.
\end{itemize}

\section*{Treatment Options}
\begin{itemize}
 \item First-line treatment is selected by severity and may combine topical benzoyl peroxide, a topical retinoid, azelaic acid, salicylic acid, or other recommended topical medicines. Topical antibiotics should not be used alone and are generally combined with benzoyl peroxide to reduce resistance.
 \item Oral doxycycline or another appropriate oral antibiotic may be used for moderate-to-severe inflammatory acne for the shortest effective duration, together with suitable topical therapy and local antimicrobial guidance. Tetracyclines are not used during pregnancy and require clinician review for contraindications and adverse effects.
 \item Combined oral contraceptives or spironolactone may help selected patients with acne linked to androgen activity, after review of contraindications, pregnancy possibility, blood pressure, drug interactions, and monitoring needs.
 \item Isotretinoin is reserved for severe, scarring, or treatment-resistant acne under specialist supervision. It is highly teratogenic and requires strict pregnancy-prevention and monitoring procedures where applicable.
 \item Intralesional corticosteroid injections may be considered by a dermatologist for selected large, painful nodules; scar treatment is considered after active acne is controlled.
\end{itemize}

\section*{Self-Care and Home Management}
\begin{itemize}
 \item Use a gentle, non-comedogenic cleanser once or twice daily and avoid harsh scrubbing, picking, or over-washing.
 \item Apply an over-the-counter benzoyl peroxide or salicylic acid product as directed; start gradually if irritation occurs. Benzoyl peroxide can bleach fabrics.
 \item Avoid picking or squeezing lesions to reduce scarring and prevent secondary infection.
 \item Use oil-free, non-comedogenic moisturisers and sunscreens labelled non-comedogenic.
 \item If pregnant, trying to become pregnant, or breastfeeding, ask a clinician before using acne medicines; topical retinoids, isotretinoin, hormonal therapy, and oral tetracyclines should be avoided or reviewed carefully. Benzoyl peroxide, azelaic acid, and limited topical salicylic acid may be options in pregnancy after professional advice.
\end{itemize}

\section*{References}
\begin{thebibliography}{99}
\bibitem{acne:ref1} American Academy of Dermatology. Acne Clinical Guideline. Updated 2024.
\bibitem{acne:ref2} National Institute for Health and Care Excellence. Acne Vulgaris: Management. NICE guideline NG198.
\bibitem{acne:ref3} American College of Obstetricians and Gynecologists. Skin Conditions During Pregnancy. Current online guidance.
\bibitem{acne:ref4} MedlinePlus. Acne. U.S. National Library of Medicine. Current online resource.
\bibitem{acne:ref5} Zaenglein AL, Pathy AL, et al. Guidelines of care for the management of acne vulgaris. \textit{J Am Acad Dermatol}. 2016;74(5):945--973.
\end{thebibliography}
